
%%%%%%%%%%%%%%%%%%%%%%%%%%%%%%%%%%%%%%%%%%%%%%%%%%%%%%%%%%%%%
\section{模型的改进与推广}

\subsection{模型的改进}

后续改进可从三个方向展开。其一是把观测网构型从"中心加同心环"推广到一般的点集优化。零漏测约束本质上是一族关于点集位置的非凸约束，其可行域的边界由"某处恰好出现 $180^\circ$ 的最大方位间隔"这一临界条件刻画。可以用逐点贪心剔除配合局部扰动的启发式在更大范围内搜索，也可以把目标域离散为细网格后用集合覆盖的整数规划给出更紧的下界，预期可在保持零漏测的前提下把 $37$ 点进一步压缩。其二是把问题二的定位精度与机动成本合并为单一目标，构造拉格朗日形式的综合目标函数，求其关于成本权重的整条最优解路径，从而把问题二与问题三的决策统一在一个框架内，直接给出精度与时间的权衡前沿。其三是引入对干扰源分布的自适应学习：在任务执行过程中根据已清除干扰源的落点不断更新空间先验，动态调整观测网的疏密，对热点区域加密、对冷区稀疏化，使策略能够适应非均匀分布的实际场景。

\subsection{模型的推广}

本文建立的模型与算法具有较强的可迁移性。问题一的半平面求交与凸多边形直径算法可推广到任意多站多目标的角度交会定位场景，包括雷达组网定位、声源追踪与移动通信中的到达角定位；其中"直径圆能否覆盖定位区域"的判据可直接用于评估以圆形置信域近似凸多边形置信域的保守程度，为传感器管理中的门限设置提供依据。问题二给出的误差传播公式与最优布站结论 $b^{*}=\sqrt2r$ 适用于一切仅有角度观测量的定位系统，其结论指出最优交会角并非 $90^\circ$ 而是 $54.74^\circ$，这一修正对无人机的多站协同测向布站、车载测向系统的机动路径规划具有直接的指导意义。问题三的"保证覆盖网加分层闭环"框架可推广到无人机群的区域搜索、污染源排查、灾后生命探测与农业植保的路径规划等需要在有限时间内保证穷尽覆盖的搜索任务，其"最少观测点数由覆盖定理给出、行程由旅行商下界刻画"的思路具有很强的通用性。问题四的方位包围判据则适用于任何存在方向性观测盲区的探测任务，例如带遮蔽的传感器网络部署、雷达盲区规避、红外探测器的朝向设计以及航天器对地观测中的太阳规避问题；其"覆盖不等于包围、观测的几何结构比观测的数量更重要"这一结论具有较普遍的工程价值。
